% ============================================================
% 片段：s01  引言
% 对应原文 PDF 第 3--5 页（印刷页码 3--5）
% ============================================================

\section{引言}
\label{sec:1}

% ==== 原文 PDF 第 3 页 ====

\begin{keytakeaways}
  \begin{itemize}[leftmargin=1.4em,itemsep=2pt,topsep=2pt]
    \item \textbf{物理图像}：应把物理理论连同其耦合空间作为一个“族”来研究；微扰与非微扰自由度在不同
          耦合区域可以互换，这种互换的物理表现就是各种对偶。镜像对称是其中最早被发现的一种：
          它把 Calabi--Yau 流形的凯勒形变与复结构形变联系起来。
    \item \textbf{数学工具}：$\mathcal{N}=(2,2)$ 超共形代数、手征环、形变空间与模空间、特殊几何、
          霍奇结构的变化、Gromov--Witten 不变量、BCOV 全纯反常方程。
    \item \textbf{关键结论}：拓扑弦理论是“只保留手征环算符”并与二维引力耦合的拓扑场论；其配分函数按
          拓扑弦耦合常数作亏格展开，见式~\eqref{eq:1.1}。B 模型因只需处理 $Z^*$ 的复结构形变而便于计算，
          A 模型则给出 $Z$ 的计数信息；高亏格振幅具有有限个生成元的多项式结构。
  \end{itemize}
\end{keytakeaways}

% [Tutor Note: 含糊点] “生成元变号给出同构”一句省略了镜像对称在 SCFT 层面的构造
\begin{yizhu}[含糊点]
原文说“某些生成元变号会给出一个同构 \cite{ref4}”，这是对镜像对称在最抽象层面的陈述：在 $\mathcal{N}=(2,2)$
超共形代数中，把 $U(1)$ 流的符号（即左右手征荷的符号）作适当翻转后，$(a,c)$ 环与 $(c,c)$ 环互换，
从而两个手征环作为代数同构。作者省略了具体构造；细节可查文献 \cite{ref4}（Dixon 关于世界面性质的综述）
以及 \cite{ref5}（Lerche--Vafa--Warner 关于 $N=2$ 手征环的经典论文）。注意这里的同构是“在手征环层面”的，
并非逐点意义上的靶空间等距。
\end{yizhu}

超对称规范理论与弦理论的研究，一直是数学与物理之间卓有成效的互动领域，双方都借此推进了知识与理解的
边界。物理方面得到的一条关键经验是：以族（family）的观点来看待物理理论非常有用，一个理论族中的成员
通过形变相互关联。这相当于把理论与它们的耦合空间放在一起研究。物理理论通常用轻自由度的拉格朗日量
及其相互作用来表述。当相互作用由小耦合主导时，微扰技术可以用来计算物理可观测量的关联函数。然而微扰论
并不考虑非微扰效应——它们是具有有限作用量的非平庸场组态，必须作为微扰论的各种背景来考虑，且一般难以
研究。当理论的耦合发生变化时，微扰论随之失效，非微扰效应变得越来越重要，这表明原来对轻自由度的选择
不再有效。因此，在耦合空间的不同区域，微扰自由度与非微扰自由度可能互换。

对微扰与非微扰自由度的相互作用、以及它们的互换以对偶形式在物理上的表现，认识上的一次大跃进来随
Seiberg 与 Witten（SW）\cite{ref1} 的开创性工作而来。该工作利用全纯性，以及有效耦合空间中由物理粒子变成
无质量所导致的奇点处的预期行为，确定了包含所有非微扰效应在内的精确低能有效理论。其关键洞见在于：
确定精确耦合的问题可以映射为确定一条曲线的周期的等价数学问题。这一辅助几何可以赋予物理意义——把它
等同于 IIB 型弦紧化中 Calabi--Yau（CY）三维流形几何的一部分 \cite{ref2}，从而把这一理解嵌入弦对偶的更大框架
（综述见 \cite{ref3}）。

% [Tutor Note: 物理图像] 周期问题与特殊几何的等价是整个讲义的出发点
\begin{yizhu}[物理图像]
Seiberg--Witten 解的核心不是“算出了一个耦合”，而是把强耦合信息编码进一个辅助几何：低能有效作用量由
一条椭圆曲线的周期 $\oint_\gamma \lambda_{\mathrm{SW}}$ 的解析延拓给出，耦合常数的强弱对应周期空间中的
不同区域。原文随后把它认作 IIB 型弦在 CY 三维流形上紧化的复结构模空间的一部分，这正是第 3、4 节
“特殊几何”的先声：模空间的特殊几何（凯勒势、Gauss--Manin 联络、汤川耦合）同时支配精确有效作用量与
拓扑弦的亏格零振幅。建议对照原文 \cite{ref1,ref2} 阅读，并留意原文在此处并未给出具体曲线。
\end{yizhu}

较早出现的一种特殊物理对偶便是镜像对称。人们观察到，在某些弦紧化背后起作用的
$\mathcal{N}=(2,2)$ 超共形代数中，把一些生成元变号会给出一个同构 \cite{ref4}。这个同构还确认了超共形代数中的
一个算符环，即手征环 \cite{ref5}。当这个代数由定义在不同靶空间 $Z$ 与 $Z^*$ 上的非线性 sigma 模型实现时，
形变问题以及用超共形代数抽象定义的同构便获得了数学意义，见 \cite{ref6} 及其中的文献。\footnote{关于镜像对称，
如今已有大量文献，特别参见 \cite{ref7,ref8,ref9,ref10} 及其中的文献，以更完整地了解导致镜像对称的发展。} 把超共形
代数表示中的态截断为由手征环中算符所创生的态，并对所有二维拓扑求和，就定义了拓扑弦理论；两个同构的
版本分别称为 A 模型与 B 模型。它们分别探测镜像 Calabi--Yau 三维流形 $Z$ 与 $Z^*$ 的凯勒形变与复结构
形变。

% ==== 原文 PDF 第 4 页 ====

B 模型更便于计算，因为它的形变是 $Z^*$ 的复结构形变，由霍奇结构的变化所刻画。镜像对称通过给出
镜像映射来建立，镜像映射是形变空间某一片区内一组特殊的局部坐标。它们给出到 A 模型的映射，因为它们
自然地与手征环相联系。

在模空间的特殊轨迹上，A 模型数据给出 CY 流形 $Z$ 的计数信息。利用构造五次超曲面镜像流形的方法 \cite{ref11}，
人们得以 \cite{ref12} 计算 B 模型中与霍奇结构变化相关的数据，并对 Gromov--Witten（GW）不变量作出预言。
高亏格 GW 不变量可以重求和，给出由 M 理论在 $Z$ 上紧化得到的五维理论中 BPS 态的整数重数 \cite{ref13,ref14}。
此外，支配形变空间的特殊几何使人能够计算预势 $F_0(t)$，它支配由 IIA（IIB）型弦在 $Z$（$Z^*$）上
紧化得到的四维理论的精确有效作用量。各种四维规范理论可以被几何工程 \cite{ref15,ref16}，而镜像对称可用来计算
它们的精确有效作用量。（另见 \cite{ref17} 及其中的文献。）

% [Tutor Note: 易错点] 五维 BPS 态计数与高亏格 GW 不变量的关系
\begin{yizhu}[易错点]
原文说“高亏格 GW 不变量可以重求和给出五维理论中 BPS 态的整数重数 \cite{ref13,ref14}”——这里的重求和是
Gopakumar--Vafa 意义下的拓扑弦自由能重排：把按亏格展开的 $F^{(g)}$ 重新组合成按 BPS 态简并数
$n^g_Q$ 的整数展开。核对时注意 (a) 五维理论由 M 理论在 CY 三维流形上紧化得到，其 BPS 粒子来自
M2 膜包裹全纯曲线；(b) 整数重数出现在 5d 的 $N=1$ 超多重态计数中，而非 4d 的 $N=2$ 理论；
(c) 原文的 $F_0(t)$ 只是亏格零部分，其导数给出四维规范耦合，不要把两者混为一谈。详见 \cite{ref13,ref14}。
\end{yizhu}

预势是拓扑弦理论的亏格零自由能；拓扑弦理论按控制高亏格展开的耦合常数做微扰定义。拓扑弦理论的配分
函数——它体现对形变空间中局部坐标的依赖——具有如下形式：
\begin{equation}
  \mathcal{Z}(t,\bar t)
  = \exp\left(\sum_g \lambda^{2g-2}\,\mathcal{F}^{(g)}(t,\bar t)\right).
  \tag{1.1}
  \label{eq:1.1}
\end{equation}

% [Tutor Note: 易错点] 亏格展开的幂次约定与 $\lambda$ 的含义
\begin{yizhu}[易错点]
此处推导极易出现符号/幂次错误：$\lambda^{2g-2}$ 中 $\lambda$ 是拓扑弦耦合常数，$g=0$ 项为
$\lambda^{-2}\mathcal{F}^{(0)}$，即预势在弱耦合下主导；$g=1$ 项与 $\lambda$ 无关。不同文献对自由能的
记号不一致（$\mathcal{F}_g$、$F_g$、$F^{(g)}$ 都有使用），本讲义在正文中混用 $F^{(g)}$ 与
$\mathcal{F}^{(g)}$，转写时照原文保留。另注意原文在讨论配分函数时写 $\mathcal{Z}$，而在讨论振幅时写
$F$，两者不要混写；建议核对原文第 4 页式 (1.1) 与第 27 页式 (4.1) 的约定。
\end{yizhu}

在文献 \cite{ref18,ref19} 中，Bershadsky、Cecotti、Ooguri 与 Vafa（BCOV）发展了高亏格拓扑弦自由能的理论与
性质，提出了递推方程——全纯反常方程，以及用费曼图求解它们的方法。对于完整的配分函数，这些方程取
热方程的形式 \cite{ref19,ref20}，并且可以被解释 \cite{ref20} 为：当配分函数被看作与 $H^3(Z^*)$ 的几何量子化相联系的
波函数时，它描述了配分函数的背景无关性。

利用形变空间的特殊几何，人们在五次超曲面及相关的单参数形变族上证明了高亏格振幅在有限个生成元中的
多项式结构 \cite{ref21}，并推广到任意靶 CY 流形 \cite{ref22}，另见 \cite{ref23,ref24,ref25}。这一多项式结构再加上适当的边界条件，
增强了高亏格振幅的可计算性。此外，这些多项式生成元有望弥合差距，帮助构造给定靶空间对偶群所对应的
恰当模形式——该对偶群由 CY 流形的特殊几何反映出来。

% ==== 原文 PDF 第 5 页 ====

本讲义旨在对上述一些思想作可读而简明的介绍。本讲义远非自足，所涵盖的题目也只是这一领域中大量进展的
一小部分；文中将给出原始文献以及许多优秀综述的指引。由于其性质，这里呈现的材料与作者自己的工作非常
接近，实际上大量依赖文献 \cite{ref22,ref26,ref27,ref28}。

% [Tutor Note: 物理图像] 高亏格振幅的多项式结构为何重要
\begin{yizhu}[物理图像]
“多项式结构”指的是：所有亏格 $g\ge 2$ 的自由能 $\mathcal{F}^{(g)}$ 都可以写成有限个生成元
（如 $S$、$S^z$、$S^{zz}$、$P$ 等）的多项式，其系数是模参数的函数。物理上，这反映了高亏格振幅的非全纯
部分完全由低亏格的“骨架”数据经反常方程递推决定，全纯部分则被边界条件（锥形流形处的行为、
大体积处的多项式行为）固定。数学上它把无穷多个亏格的计算化归为有限个生成元的代数运算，并提示
这些生成元应当组织成对偶群的模形式。详细构造见第 4.4--4.5 节。
\end{yizhu}

本讲义的大纲如下。第 2 节回顾 $\mathcal N=2$ 超共形代数及其手征环与形变，并讨论如何由形变算符张成的
手征环子环构造态丛；此外描述 $\mathcal{N}=(2,2)$ 超共形代数作为非线性 sigma 模型的几何实现。第 3 节
概述如何先把物理可观测对象截断到手征环、再与引力耦合而得到拓扑弦理论；讨论两种不等价的方式，即
A 模型与 B 模型，以及手征环在两侧的几何解释；回顾特殊几何及其在霍奇结构变化中的 B 模型实现；并以
五次超曲面及其镜像为例，说明如何构造镜像几何并求出镜像映射。第 4 节回顾 BCOV 反常方程及其递推解；
利用特殊几何可以证明高亏格振幅的多项式结构，并在实现物理预期的边界条件时用于更简便的计算；最后给出
把该结构应用于五次超曲面的例子。
